
% \subsection{Описание функция и требования к блоку экспорта и анализа}

% НУЖНО НАПИСАТЬ!!!

% После выполнения необходимых расчетов выполняется совмещение расчитанных показателей с полученными критериями риска для объекта, после чего происходит их интерполяция и интерпретация, продолжающаяся процедурами экспорта и сериализации полученных результатов.

% НУЖНО НАПИСАТЬ!!!

\subsection{Описание функций и требования к блоку экспорта и анализа}

После выполнения необходимых расчетов выполняется совмещение вычисленных
волновых показателей, описанных в разделе~\ref{sec:wer_calculation} на стр.~\pageref{sec:wer_calculation} с
полученными критериями риска для исследуемого объекта береговой линии.
На данном этапе каждому трансекту или сегменту линии присваивается итоговый
ранг волнового воздействия $\mathrm{WER}$ в соответствии с таблицей
рангов~\ref{tab:cwef_percentiles}, после чего формируется совокупная карта 
оценки уязвимости побережья.

Интерполяция рассчитанных показателей применяется с целью обеспечения
пространственной непрерывности результатов вдоль береговой линии.
В случае неравномерного расположения исходных трансектов или наличия
пространственных пробелов в метеорологических или батиметрических данных
значения $\mathrm{CWEF}$ и $\mathrm{WER}$ (показатели выведены в разделах~\ref{eq:CWEF} и~\ref{sec:wer_calculation}) они восстанавливаются методом
линейной интерполяции.
Интерпретация результатов предполагает отнесение каждого участка береговой линии к одному из пяти классов волновой экспозиции в соответствии с выработанными пороговыми значениями (раздел~\ref{sec:wer_usage} на стр.~\pageref{sec:wer_usage}).

Завершающим этапом работы блока является экспорт и сериализация
полученных результатов.
Блок обеспечивает сохранение пространственных данных в форматах
\texttt{GeoJSON} и \texttt{GeoPackage}, пригодных для
последующей визуализации и анализа в ГИС-средах, в том числе
\textit{QGIS}.
Помимо векторных геопространственных форматов, предусмотрен экспорт
табличных данных в формате \texttt{CSV}/\texttt{XLSX} для статистической
обработки и отчетности, а также вывод картографических материалов в
форматах, совместимых с системами автоматизированного проектирования (CAD).

% К блоку экспорта и анализа предъявляются следующие функциональные
% требования:

% \begin{enumerate}
%     \item \textbf{Корректное совмещение показателей.}
%           Все рассчитанные волновые параметры должны быть
%           пространственно привязаны к единой системе координат объекта,
%           а присвоение рангов $\mathrm{WER}$ -- осуществляться строго
%           по формуле~\eqref{eq:wer} без потери атрибутивных данных.

%     \item \textbf{Пространственная интерполяция.}
%           При наличии пробелов в данных блок должен выполнять
%           интерполяцию значений вдоль береговой линии с возможностью
%           задания метода интерполяции (линейный, сплайн).

%     \item \textbf{Многоформатный экспорт.}
%           Результаты должны экспортироваться не менее чем в три
%           формата: \texttt{GeoJSON}, \texttt{GeoPackage} и
%           \texttt{CSV}; при этом метаданные (система координат,
%           дата расчета, версия программного обеспечения) должны
%           сохраняться в атрибутах файла.

%     \item \textbf{Человекочитаемая интерпретация.}
%           В выходных данных наряду с числовыми значениями
%           $R_i$ и $\mathrm{WER}$ должны присутствовать текстовые
%           метки классов экспозиции (например, «низкая»,
%           «умеренная», «высокая») согласно принятой классификации.

%     \item \textbf{Интеграция с визуализационным блоком.}
%           Блок должен передавать геопространственные данные во
%           встроенный модуль визуализации для автоматического
%           построения тематических карт волновой экспозиции
%           и графиков временны́х рядов $\mathrm{CWEF}(t)$.
% \end{enumerate}